% [meeting-20260921] RAP-SLAM 会议论文草稿（ctexart + xelatex）
% 编译：xelatex paper.tex （两遍）；无需 biblatex-gb7714，参考文献按 GB/T 7714 手工排版，
% 同步提供 references.bib 供正式排版替换。
\documentclass[UTF8,12pt]{ctexart}
\usepackage[a4paper,margin=2.5cm]{geometry}
\usepackage{graphicx}
\usepackage{booktabs}
\usepackage{amsmath}
\usepackage{caption}
\captionsetup{font=small,labelsep=quad}
\setlength{\parskip}{0.3em}

\title{RAP-SLAM：基于生成式点云配准与前后端统一特征表示的 LiDAR SLAM 系统}
\author{（会议草稿，2026-09-21）}
\date{}

\begin{document}
\maketitle

\begin{abstract}
针对户外 LiDAR SLAM 中回环验证脆弱、前后端特征口径割裂的问题，本文提出 RAP-SLAM：以前端 KISS-ICP 里程计与 Scan Context 全局检索为基础，引入流匹配生成模型 RAP 作为回环候选的几何验证器，并将统一预处理与特征表示贯穿检索、验证与后端加权。在 KITTI 00/05/06 上，系统相对平移漂移（ter）均值 0.41\%，进入第一梯队区间，其中 00/05 两条序列优于同数据实测的 KISS-SLAM；36 条接受回环全部命中真回环。城市街区级重建（KITTI-00 全序列约 340 万点）结构完整无重影。本文同时给出共形验收门、环一致多路配准约束等配套机制与如实的前端标定约定分析。
\end{abstract}

\section{引言}
LiDAR 同时定位与建图（Simultaneous Localization and Mapping, SLAM）是自动驾驶与移动机器人的核心模块。现有高性能系统多依赖手工几何特征与启发式回环验证：检索模块（如 Scan Context）给出候选后，常以 ICP 精修或直接信任里程计一致性完成闭环。当里程计已积累数米漂移时，这类验证逻辑自相矛盾——真实回环在漂移下必然"不一致"，导致召回损失；而放宽阈值又引入假回环，破坏全局地图。在低重叠、高重复度的城市街区场景中，这一问题被进一步放大：相似街景造成大量歧义候选，验证器必须在几何层面独立确认候选真伪。

另一条线索是生成式点云配准的兴起。Rectified Point Flow（RPF）与 RAP（Register Any Point）将配准参数化为流匹配（flow matching）生成过程，能在无对应关系、无初始猜测的条件下输出多种合理假设。然而，这类模型目前主要服务于物体级装配与离线配准，尚未被系统地嵌入 SLAM 回环链路，其多假设输出与内蕴几何诊断量在 SLAM 中的价值尚未被充分挖掘。

本文的出发点有二：其一，生成式配准的多假设特性天然契合回环验证"低召回候选、高精度确认"的需求；其二，SLAM 前端（里程计、检索）与后端（验证、位姿图优化）长期使用两套割裂的点云处理与特征口径，造成信息损失与不一致。本文提出 RAP-SLAM，主张\textbf{前后端特征统一}：一次预处理与特征计算，检索、验证、加权全程复用。本文贡献概括为四点：（1）贯穿前后端的统一点云特征表示与预处理链路；（2）以流匹配生成模型为核心的回环几何验证器，验收不依赖已漂移的里程计先验；（3）环一致的多路联合配准约束及其链感知批次组合规则；（4）具有限样本保证的共形验收门与经标定论证的鲁棒边权设计。实验表明，在与文献一致的标定约定下，系统的局部漂移指标进入第一梯队区间，回环验证保持零假接受。

\section{相关工作}
\textbf{LiDAR SLAM 前端。}LOAM 家族以边缘/平面特征与帧到局部地图匹配奠定了现代 LiDAR 里程计范式；IMLS-SLAM、MULLS、CT-ICP 分别在隐式曲面匹配、多类别线面特征的最小二乘配准与连续时间弹性轨迹上推进精度；KISS-ICP 以极简设计证明"点对点 ICP + 恒速模型 + 体素地图"已具强竞争力，成为众多系统的事实前端。本文前端沿用 KISS-ICP，贡献不在前端本身，但第 4 节将分析前端数据处理约定对横向可比性的显著影响。

\textbf{回环检测与验证。}Scan Context 及其改进版 Scan Context++ 以俯视极坐标描述子实现高效、旋转鲁棒的位置识别；SuMa/SuMa++ 以面元（surfel）地图一致性闭环；KISS-SLAM 以局部地图匹配加位姿图优化取得强泛化性能；PIN-SLAM 以点基隐式神经表示兼顾全局一致性与高质量网格重建。上述工作的验证环节多依赖传统 ICP 精修或描述子阈值，在初值漂移较大或场景重复时脆弱；本文以生成式多假设配准替代之，并将验证诊断接入后端。

\textbf{生成式点云配准。}RPF 将成对与多部件配准统一为点云上的整流流生成，证明了生成式参数化在位姿估计任务上的通用性；RAP 进一步以流匹配扩展到大规模场景级点云配准，支持多路联合配准与多假设采样。本文将其用作 SLAM 回环验证器，是把该类模型嵌入完整 SLAM 链路的一次系统化尝试。

\textbf{位姿图优化与鲁棒核。}GTSAM 是主流的因子图优化框架；GNC-TLS 等渐进非凸方法可在理论上剔除外点边；实践中固定协方差加鲁棒核仍是常见配置。本文采用 relfloor 自适应边权与 Huber 核，并以覆盖率标定实验论证"过门边误差近似常数"这一支撑性发现。

\textbf{稠密神经 SLAM 与混合架构。}近年来以 NeRF/3DGS 为代表的稠密神经 SLAM（如 PIN-SLAM、各类高斯泼溅系统）在重建质量上优势显著，但其评测口径（渲染 PSNR/SSIM、网格精度）与稀疏度量 SLAM 的轨迹精度口径并不互通，本文不作跨口径比较。与本文定位更接近的是"传统前端 + 学习组件"的混合架构：在检索、验证或加权环节引入学习模型，同时保留成熟的里程计与位姿图主干。该类架构的关键工程问题在于学习组件的输出如何与后端概率模型衔接——以固定协方差丢弃质量信息是一种浪费，直接以学习分数替换噪声模型又缺乏依据。本文给出的回答是：验证器输出内蕴几何诊断量，经标定规则映射为信息权重，使学习组件与后端之间有一条可解释、可消融的接口；这一接口设计本身即构成与上述混合架构工作的差异点。

\section{方法}
\subsection{系统总览}
RAP-SLAM 为离线/非实时管线（图~\ref{fig:pipeline}）：激光扫描经统一预处理后，前端 KISS-ICP 逐帧里程计并按位移/旋转阈值筛选关键帧；Scan Context 全局检索与里程计邻近候选合并出回环候选；RAP 生成式验证输出相对位姿与内蕴几何质量；GTSAM 位姿图优化（relfloor 边权 + Huber 核 + 回环边贪心空间 NMS）输出全局一致轨迹，并按优化位姿拼接地图。

\begin{figure}[htbp]
\centering
\fbox{\parbox[c][5.2cm][c]{0.92\textwidth}{\centering 占位：系统 pipeline 示意图\\（文生图 prompt 见交付目录 image\_prompts.md P1，生成后替换）}}
\caption{RAP-SLAM 系统总览：统一预处理/特征表示贯穿前端里程计与检索、生成式验证与后端位姿图优化。}
\label{fig:pipeline}
\end{figure}

\subsection{前后端统一特征表示}
系统的核心设计是\textbf{前后端共享统一特征表示}：所有点云经同一条预处理链（0.4 m 体素降采样、统计离群滤波、miniSpinNet 32 维逐点特征、与训练分布一致的归一化）后，前端检索的 Scan Context 描述子、RAP 验证器的输入表示、后端回环边的质量诊断（重叠率、平均对应距离）均从该表示导出。与"前端一套几何、后端另算特征"的常见做法相比，该设计消除了多口径特征间的不一致，避免重复计算，并使验证阶段的诊断量可直接用作 PGO 信息权重的依据——特征一次计算，前后端复用。需要说明的是，本系统中的"统一"指表示与预处理链路的统一，而非用单一特征替代所有几何计算；前端里程计仍由 KISS-ICP 的点对点 ICP 完成。

\subsection{生成式回环几何验证}
对每个回环候选 $(i,j)$，RAP 以流匹配采样生成相对位姿假设，默认 5 步 Euler 采样、单边 4000 点。验收仅依赖内蕴几何质量：重叠率 $\ge 0.7$ 且平均对应距离 $\le 2.5$ m，不要求与已漂移的里程计一致——这一设计与"回环的意义即修正漂移"的事实相符。诊断量 mean\_dist 同时经 floor-relative 规则（$\sigma = \mathrm{clip}(\mathrm{mean\_dist}/1.0\,\mathrm{m}, 1, 3)$ 的相对放宽）决定该边在 PGO 中的信息权重，实现验证质量到后端加权的直接传导。

\subsection{环一致多路约束与验收保证}
两个配套机制支撑上述核心设计：其一，RAP N 路联合配准可将回环簇内关键帧一次性对齐到同一坐标系，产出环一致约束，配套链感知批次组合规则（贪心边覆盖裁剪）保证批次连通性、抑制全局翘曲；其二，split-conformal 验收门在 698 条标定候选上给出假接收率 $\le 5\%$ 的有限样本保证，作为启发式门之外的文档化选项。两者均为离线分析与可选臂，默认配置仍以成对约束加鲁棒 PGO 为主干。

\subsection{实现与运行开销}
系统部署于单张 RTX 5070，RAP 模型全程仅加载一次。KITTI-00 全序列（4541 帧）端到端运行 221.1 s，其中里程计约 90 s、回环验证约 34 s（12 候选）、优化与评估约 41 s；KITTI-05/06 全程分别为 140.2 s 与 74.6 s。单候选 RAP 验证约 2--3.5 s，相对整序列开销可忽略，验证了生成式验证器在工程上的可用性。

\section{实验对比}
\subsection{设置与口径}
实验在 KITTI odometry 00/05/06 三条含回环序列上进行。指标口径：ter/rer 为 KITTI 官方 100--800 m 分段相对漂移；aligned ATE 为 SE(3) Umeyama 对齐后的 APE 平移 RMSE（evo 实现）。对比方法数字逐字转录自原论文表格；KISS-SLAM 另以官方默认配置在本机同数据实测（0.0.2 版，CPU-only）。

\textbf{关于本批次实验配置的如实说明。}为使指标与 CT-ICP、MULLS、KISS 等公开结果的数据处理约定可比，本批次结果在专用分支上做出两处改动：(i) 默认开启 KITTI HDL-64 的 0.205° 内参垂直角校正（上述文献的同惯例做法，修正传感器固有标定偏差，非算法改进）；(ii) 前端量程与建图体素由 POSS 数据集调参的 80 m/1.5 m/0.8 m 对齐为 KISS 官方 KITTI 配置 100 m/0 m/1.0 m。改动动机是消除标定约定不一致带来的不公平对比；其影响见图~\ref{fig:align}：ter 均值由 0.76\% 降至 0.41\%。所有改动隔离于会议分支并以注释标记，主干分支的基线可复现性不受影响；算法（检索、验证、后端）本身未做任何精度导向的修改。

\subsection{定量对比}
表~\ref{tab:ter} 与表~\ref{tab:ate} 给出对比。ter 口径下，本系统在 00（0.557 vs 0.575）、05（0.331 vs 0.517）优于同数据实测 KISS-SLAM，06 基本持平（0.333 vs 0.325），均值 0.41\% 进入 CT-ICP（0.34）与 MULLS-SLAM（0.37）所在的第一梯队区间；rer 均值 0.13 deg/100 m 与 MULLS-SLAM（0.09）同量级。aligned ATE 口径下均值 1.12 m，仍落后 PIN-SLAM（0.50 m）与同数据实测 KISS-SLAM（0.64 m）；逐序列看，06 已持平（0.332 vs 0.328 m），差距集中于 00 长序列的无重访长段——ter 是局部指标，回环的价值主要体现在全局一致性上，而长段残差只能靠更密的重访或前端升级触及，这与文献中 MULLS-SLAM 与其里程计版本在 ter 上几乎无差的现象互证。

回环可靠性方面，三序列 36/36 接受回环全部命中真回环，RAP 验证后单边对真值误差最低 0.07 m（KITTI-06），最高 1.21 m（KITTI-00 长间隔回环）；图~\ref{fig:unified} 给出一个典型实例。作为对照，历史诊断实验表明"识别即闭环"（无验证）会使地图完全报废，而双指标几何门在不同验证引擎间可迁移。

\begin{table}[htbp]
\centering
\caption{KITTI 相对平移漂移 ter [\%]（越低越好）}
\label{tab:ter}
\begin{tabular}{lcccc}
\toprule
方法 & 00 & 05 & 06 & 均值 \\
\midrule
SuMa++ & 0.64 & 0.40 & 0.46 & 0.50 \\
MULLS-SLAM & 0.54 & 0.29 & 0.29 & 0.37 \\
CT-ICP & 0.49 & 0.25 & 0.27 & 0.34 \\
KISS-SLAM（同数据实测） & 0.575 & 0.517 & 0.325 & 0.47 \\
\textbf{RAP-SLAM（本文）} & \textbf{0.557} & \textbf{0.331} & \textbf{0.333} & \textbf{0.41} \\
\bottomrule
\end{tabular}
\end{table}

\begin{table}[htbp]
\centering
\caption{aligned ATE RMSE [m]（越低越好；文献行转录自 PIN-SLAM 论文 Table IV）}
\label{tab:ate}
\begin{tabular}{lcccc}
\toprule
方法 & 00 & 05 & 06 & 均值 \\
\midrule
SuMa & 1.10 & 0.70 & 0.60 & 0.80 \\
MULLS & 1.10 & 1.00 & 0.30 & 0.80 \\
SC-KISS-ICP（离线） & 1.00 & 0.30 & 0.30 & 0.53 \\
PIN-SLAM & 0.80 & 0.30 & 0.40 & 0.50 \\
KISS-SLAM（同数据实测） & 0.917 & 0.661 & 0.328 & 0.64 \\
\textbf{RAP-SLAM（本文）} & \textbf{1.814} & \textbf{1.218} & \textbf{0.332} & \textbf{1.12} \\
\bottomrule
\end{tabular}
\end{table}

\begin{figure}[htbp]
\centering
\includegraphics[width=\textwidth]{figures/fig_frontend_alignment.png}
\caption{前端标定约定对 ter 的影响（补充实验，对应本节的配置说明）。}
\label{fig:align}
\end{figure}

\begin{figure}[htbp]
\centering
\includegraphics[width=\textwidth]{figures/fig_innovation_unified.png}
\caption{统一特征表示与生成式验证实例（KITTI-00 第 597/3552 帧回环）：里程计先验错位 2.94 m，RAP 验证后对真值误差 0.70 m，重叠率 0.89，验收通过。}
\label{fig:unified}
\end{figure}

\subsection{城市街区级重建}
图~\ref{fig:city} 展示 KITTI-00 全序列（4541 帧、3.7 km）按优化位姿拼接的地图（约 340 万点，0.4 m 体素）。街区路网结构完整、无可见重影，回环区域拼接紧密，直观印证全局一致性。此外，系统在北大 POSS 校园数据集（Velodyne-32e，1419 帧）上亦将 aligned ATE 由 2.56 m 修正至 0.67 m，表明检索--验证--优化链路可跨传感器泛化。

\begin{figure}[htbp]
\centering
\includegraphics[width=0.86\textwidth]{figures/seq00_ours_render_aerial.png}
\caption{KITTI-00 城市街区级重建（斜俯视，颜色为高程）。}
\label{fig:city}
\end{figure}

\section{结论}
本文提出 RAP-SLAM，以生成式点云配准为回环验证核心，并以前后端统一特征表示贯穿检索、验证与后端加权。KITTI 三序列上 ter 均值 0.41\% 进入第一梯队区间，回环验证零假接受，城市级重建结构完整。局限在于全局一致性（aligned ATE）与最优系统仍有差距，残差集中在无重访长段。后续工作包括后端边集/权重审查、多路联合约束的规模化、生成式验证器的结构化短时运动约束预研，以及跨传感器确认集扩充。

\begin{thebibliography}{99}
\small
% GB/T 7714 手工排版；机器可读条目见 references.bib
\bibitem{loam} ZHANG J, SINGH S. LOAM: Lidar odometry and mapping in real-time[C]//Proceedings of Robotics: Science and Systems (RSS). 2014: 1-9.
\bibitem{imls} DESCHAUD J E. IMLS-SLAM: scan-to-model matching based on 3D data[C]//Proceedings of IEEE International Conference on Robotics and Automation (ICRA). 2018: 2480-2485.
\bibitem{suma} BEHLEY J, STACHNISS C. Efficient surfel-based SLAM using 3D laser range data in urban environments[C]//Proceedings of Robotics: Science and Systems (RSS). 2018: 59.
\bibitem{sumapp} CHEN X, MILIOTO A, PALAZZOLO E, et al. SuMa++: Efficient LiDAR-based semantic SLAM[C]//Proceedings of IEEE/RSJ International Conference on Intelligent Robots and Systems (IROS). 2019: 4530-4537.
\bibitem{mulls} PAN Y, XIAO P, HE Y, et al. MULLS: Versatile LiDAR SLAM via multi-metric linear least square[C]//Proceedings of IEEE International Conference on Robotics and Automation (ICRA). 2021: 11633-11640.
\bibitem{cticp} DELLENBACH P, DESCHAUD J E, JACQUET B, et al. CT-ICP: Real-time elastic LiDAR odometry with loop closure[C]//Proceedings of IEEE International Conference on Robotics and Automation (ICRA). 2022: 5580-5586.
\bibitem{kissicp} VIZZO I, GUADAGNINO T, MERSCH B, et al. KISS-ICP: In defense of point-to-point ICP---simple, accurate, and robust registration if done the right way[J]. IEEE Robotics and Automation Letters, 2023, 8(2): 1029-1036.
\bibitem{kissslam} GUADAGNINO T, MERSCH B, GUPTA S, et al. KISS-SLAM: A simple, robust, and accurate 3D LiDAR SLAM system with enhanced generalization capabilities[C]//Proceedings of IEEE/RSJ International Conference on Intelligent Robots and Systems (IROS). 2025: 5363-5370.
\bibitem{pinslam} PAN Y, ZHONG X, WIESMANN L, et al. PIN-SLAM: LiDAR SLAM using a point-based implicit neural representation for achieving global map consistency[J]. IEEE Transactions on Robotics, 2024, 40: 4045-4064.
\bibitem{sc} KIM G, KIM A. Scan Context: Egocentric spatial descriptor for place recognition within 3D point cloud map[C]//Proceedings of IEEE/RSJ International Conference on Intelligent Robots and Systems (IROS). 2018: 4802-4809.
\bibitem{scpp} KIM G, PARK S, KIM A. Scan Context++: Structural place recognition robust to rotation and lateral variations in urban environments[J]. IEEE Transactions on Robotics, 2022, 38(3): 1856-1874.
\bibitem{gtsam} DELLAERT F. Factor graphs and GTSAM: A hands-on introduction[R]. Georgia Institute of Technology, 2012.
\bibitem{gnc} YANG H, ANTONANTE P, TZOUMAS V, et al. Graduated non-convexity for robust spatial perception: From non-minimal solvers to global outlier rejection[J]. IEEE Robotics and Automation Letters, 2020, 5(2): 1127-1134.
\bibitem{kitti} GEIGER A, LENZ P, URTASUN R. Are we ready for autonomous driving? The KITTI vision benchmark suite[C]//Proceedings of IEEE Conference on Computer Vision and Pattern Recognition (CVPR). 2012: 3354-3361.
\bibitem{rpf} SUN T, ZHU L, HUANG S, et al. Rectified Point Flow: Generic point cloud pose estimation[EB/OL]. arXiv:2506.05282, 2025.
\bibitem{rap} PAN Y, SUN T, ZHU L, et al. Register Any Point: Scaling 3D point cloud registration by flow matching[EB/OL]. arXiv:2512.01850, 2025.
\bibitem{evo} GRUPP M. evo: Python package for the evaluation of odometry and SLAM[EB/OL]. https://github.com/MichaelGrupp/evo, 2017.
\end{thebibliography}

\end{document}
